\section{Descripción de las Decisiones de Diseño}

\subsection{Esquema de Direccionamiento}

El bloque \texttt{172.16.192.0/19} ofrece 8.190 direcciones utilizables, lo
que permite derivar subredes de tamaño variable mediante VLSM para ajustarse
con precisión a las necesidades de cada segmento. Se optó por este enfoque
en lugar de un esquema de subredes de tamaño fijo, dado que las diferencias
de escala entre las áreas son significativas: LAN 1 requiere más de 900 hosts,
mientras que los enlaces punto a punto entre routers solo necesitan dos
direcciones. Emplear subredes de tamaño uniforme desperdiciaría una fracción
considerable del espacio de direcciones disponible.

Para IPv6, el prefijo \texttt{2001:dbe:bef9::/48} permite asignar subredes
\texttt{/64} a cada segmento de la topología sin restricción práctica de
cantidad. Se adoptó el prefijo \texttt{/64} de forma uniforme para todos los
segmentos, dado que es el tamaño mínimo requerido por SLAAC para la
autoconfiguración de direcciones en los hosts de usuario final, y simplifica
la gestión del plan de direccionamiento al eliminar la necesidad de VLSM en
la familia IPv6.

\subsection{Estrategia de Asignación de Direcciones}

Para los dispositivos de telecomunicaciones y los servidores de servicios se
adoptó asignación de direcciones estática en ambas familias de protocolos.
Esta decisión responde a que dichos dispositivos son puntos críticos de la
infraestructura: una dirección que cambie en un router o en un servidor
invalida rutas configuradas, entradas DNS y reglas de firewall, introduciendo
interrupciones difíciles de diagnosticar en producción. La asignación estática
elimina esta fuente de inestabilidad y facilita la trazabilidad durante la
administración remota.

Para los equipos de usuario final se adoptó asignación dinámica. En IPv4 se
emplea DHCP, de modo que un servidor centralizado distribuye direcciones,
máscara, puerta de enlace y servidor DNS a todos los hosts de cada LAN sin
intervención manual. En IPv6 se emplea SLAAC, mecanismo mediante el cual
cada host construye su propia dirección a partir del prefijo \texttt{/64}
anunciado por el router en sus mensajes Router Advertisement, sin requerir
un servidor dedicado. Esta combinación reduce la carga administrativa en redes
con cientos de usuarios y es coherente con las prácticas estándar en redes
empresariales de escala media.

\subsection{Protocolo de Enrutamiento}

Se seleccionó OSPF como protocolo de enrutamiento dinámico para IPv4, y su
extensión OSPFv3 para IPv6. La elección de OSPF sobre alternativas como RIP
se fundamenta en tres razones principales. Primero, OSPF es un protocolo de
estado de enlace que converge significativamente más rápido que los protocolos
de vector de distancia ante cambios en la topología, lo cual es relevante en
una red multi-sede donde una caída de enlace no debe prolongarse en
inconsistencias de ruteo. Segundo, OSPF no tiene la limitación de 15 saltos
de RIP, lo que lo hace apropiado para topologías que puedan crecer en número
de sedes. Tercero, OSPFv3 extiende de forma nativa el mismo mecanismo a IPv6,
permitiendo mantener una única lógica de enrutamiento para ambas familias de
protocolos.

Toda la topología se configura dentro de una única área OSPF (área 0), dado
que el número de routers no justifica la segmentación en áreas adicionales y
una sola área simplifica la configuración y el diagnóstico.

\subsection{Seguridad en Dispositivos de Telecomunicaciones}

El acceso local a los dispositivos se protege mediante contraseñas cifradas
en consola y en las líneas de terminal virtual. Para la administración remota
se habilita exclusivamente SSHv2, deshabilitando Telnet en todas las líneas
VTY. SSHv2 cifra la sesión completa, incluyendo credenciales y comandos,
lo que impide que un atacante con acceso pasivo a la red pueda capturar
información de configuración sensible.

\subsection{Seguridad en Servidores de Servicios}

En el servidor Linux que aloja los servicios de aplicación se aplica una
política de firewall restrictiva mediante UFW, abriendo únicamente los puertos
correspondientes a cada servicio: 80 para HTTP, 25 y 143 para SMTP e IMAP
respectivamente, y el puerto definido por la implementación del servicio de
chat. El acceso administrativo al sistema operativo se realiza por SSHv2 con
autenticación mediante clave pública, deshabilitando la autenticación por
contraseña en el demonio SSH para reducir la exposición a ataques de fuerza
bruta.

\subsection{Servicios de Aplicación}

El servicio web se implementa sobre Nginx, servidor que ofrece bajo consumo
de recursos y soporte nativo para IPv6 mediante la directiva
\texttt{listen [::]:80}, permitiendo atender conexiones de ambas familias de
protocolos sobre un único proceso. El servicio de correo se implementa con
el stack Postfix y Dovecot: Postfix actúa como agente de transferencia de
correo gestionando los procesos SMTP de envío y recepción, mientras que
Dovecot expone los buzones de usuario mediante IMAP, protocolo que permite
a los clientes acceder al correo sin descargarlo ni eliminarlo del servidor,
lo que es adecuado para un entorno corporativo con múltiples dispositivos
por usuario.

El servicio de mensajería instantánea corresponde a la implementación
desarrollada por el grupo en el reto previo, construida en C sobre sockets
TCP. El servidor opera en Linux y acepta conexiones de clientes Windows a
través de la red. Soporta múltiples usuarios simultáneos, salas de chat
con interacción en canal común, listado de usuarios conectados mediante el
comando \texttt{/list} e historial local de mensajes exportable a archivo
de texto. Para cumplir el requisito de doble stack, el socket del servidor
se configuró con la familia \texttt{AF\_INET6} y la opción
\texttt{IPV6\_V6ONLY} deshabilitada, lo que permite aceptar conexiones tanto
IPv4 como IPv6 sobre un único socket sin modificar la lógica de manejo de
clientes.